\documentclass[11pt]{article}
\usepackage[review]{acl}
\usepackage{times}
\usepackage{latexsym}
\usepackage[T1]{fontenc}
\usepackage[utf8]{inputenc}
\usepackage{microtype}
\title{Introduction Draft: Personalized Federated Steering}
\author{Anonymous submission}
\begin{document}
\maketitle
\section{Introduction}
\label{sec:introduction}

The same input can call for different kinds of outputs as a user's needs change.
An editor, for example, may want a summary that closely preserves an article's wording for one publishing format and a more extensively rewritten version for another.
Both should convey the relevant information faithfully, and useful choices may lie between these endpoints.
\emph{Controlled generation} aims to give users the ability to specify such output properties while maintaining task quality.
Steering offers one approach: an intervention along a learned behavioral direction is applied with adjustable strength, allowing a model to support different settings without being retrained for each request \citep{rimsky-etal-2024-steering,fierro-roger-2026-weight}.
The goal is a predictable relationship between the requested attribute level and the generated behavior.
For organizations that adapt language models to their own data, this control must also operate alongside their particular writing and format conventions.

Learning such control is difficult when an organization's examples cover only a limited portion of the desired behavioral range.
A publication may have thousands of summaries yet relatively few examples at the degree of rewriting an editor now requires.
Collecting more examples from the same distribution need not resolve that imbalance.
Newsroom illustrates the diversity that can exist across publications: its summaries use different mixtures of extractive and abstractive strategies, reusing source words and phrases to varying degrees \citep{grusky-etal-2018-newsroom}.
When organizations have complementary examples, their collective data offers supervision that an individual organization lacks.
This motivates learning control collaboratively in a setting where each organization, or \emph{client}, keeps its training examples local.
The opportunity arises from differences in behavioral coverage as well as differences in data quantity.

Collaboration introduces a tension between local specialization and global coverage.
Independent training can fit a client's generation conventions, but supervision for its control mechanism remains limited to its own examples.
Training a single shared model draws on a broader distribution, yet requires the same adaptation to accommodate different clients' conventions.
Personalized federated learning provides a way to combine shared knowledge with local specialization through shared and client-specific model components \citep{collins-etal-2021-fedrep,guo-etal-2025-fedsa}.
For controlled generation, however, this division must do more than support each client's typical output: the shared component must enable deliberate variation within its personalized model.
Prior personalized steering already studies behavioral intensity and transfer across models and adapters \citep{cao-etal-2024-bipo}.
We investigate the complementary question of how clients with uneven behavioral coverage can \emph{jointly learn} a control mechanism that works with private adaptation and remains useful to clients outside the training federation.

Three challenges arise in this setting.
First, the shared direction must expand a client's behavioral range while remaining compatible with its private adaptation.
That adaptation can learn the client's typical attribute level together with its writing conventions, making it difficult to change one while retaining the other.
A direction learned independently of the personalized models need not interact with them as intended.
Second, learning a direction does not determine how strongly to apply it for a requested attribute level.
The generated behavior need not respond linearly to intervention strength, so a simple scaling rule may provide uneven control across the requested range.
Third, success on participating clients may depend on the particular private adaptations encountered during training.
For collaboration to provide reusable control, a new client should be able to benefit from the learned mechanism with limited local data, rather than having to learn its own direction from the beginning.

We address these challenges through personalized federated training of a shared steering direction together with private client adaptation.
Each client learns to generate its training examples with steering active at the examples' observed attribute levels.
Clients periodically aggregate updates to the shared control mechanism while retaining their private adaptation locally.
The direction is therefore optimized in the presence of the personalized models on which it must act, drawing on behavioral examples distributed across the federation.
This joint training protocol is designed to make broader coverage available within each client's model while preserving the capacity to learn local generation characteristics.
It also makes the interaction between shared control and private adaptation an explicit part of learning.

Alongside the direction, we learn a shared monotone nonlinear scaling function that maps the requested attribute level to an intervention strength.
The direction specifies a way to change behavior, while this function determines how strongly to apply that change at different requested levels.
Its flexibility accommodates a potentially nonlinear response without prescribing a fixed proportional relationship between the request and intervention strength.
Monotonicity orders the intervention strengths; whether the resulting outputs follow the requested levels remains an empirical question.
The direction, scaling function, and private adaptations are learned jointly.
After federation, the shared control mechanism can be frozen and reused while a new client learns its own private adaptation.

% Evidence boundary: nonlinear scaling is a design choice; its advantage over
% fixed-gain or learned-gain linear scaling awaits the ablation in evaluation RQ4.
% Do not add a settled participant-superiority claim before the local gain-cap
% check recorded in experiment-log entry "Results: C1 at 8B" is resolved.
Our contributions are fourfold:
\begin{itemize}
    \item We formulate collaborative learning of continuous behavioral control under uneven client attribute coverage, with local specialization and usefulness beyond participating clients as explicit requirements.
    \item We develop a personalized federated training protocol that jointly learns a shared steering direction and scaling function together with private client adaptation.
    \item In the evaluated low-data settings on four held-out Newsroom clients, reusing the frozen control mechanism yields lower control error than learning a direction locally at the same local data and step budgets, supporting its usefulness beyond the training federation.
    \item Through publication-partitioned summarization experiments, we characterize the trade-offs among behavioral coverage, control accuracy, generation quality, and client-specific characteristics. The measured publication characteristics remain close to those of locally trained models, while comparisons with shared adaptation expose a trade-off between specialization and out-of-support control.
\end{itemize}


\bibliography{custom}
\end{document}
